% =========================================================================
% CAPITOLO 8: DECOMPOSIZIONE STAGIONALE, STAZIONARIETA E IDENTIFICAZIONE
% (Da completare nello Step 3)
% =========================================================================
\section{Decomposizione stagionale, stazionarieta e analisi delle autocorrelazioni}
\label{sec:stl_stazionarieta}

Il presente capitolo è dedicato all'analisi della struttura temporale della serie mediante decomposizione STL con periodo stagionale annuale ($s = 12$). Vengono illustrati i risultati formali dei test di radice unitaria (Augmented Dickey-Fuller e KPSS) per verificare la non stazionarietà della serie originaria e la sua riconduzione a stazionarietà mediante differenziazione prima. L'analisi è completata dall'ispezione dei correlogrammi di autocorrelazione semplice (ACF) e parziale (PACF) per la selezione degli ordini autoregressivi e di media mobile.

